% Generated from project outputs. Do not edit by hand.
\begin{table}[!htbp]
\centering
\begin{threeparttable}
\caption{Capacity-curve sensitivity: bandwidth 2016-2019 bandwidth}
\label{tab:robustness-bandwidth-2016-2019-bandwidth}
\small
\setlength{\tabcolsep}{4pt}
\renewcommand{\arraystretch}{1.15}
\begin{tabularx}{\linewidth}{@{}Xrrrrr@{}}
\toprule
Variant & Span & Q500 & Q1500 & Q3000 & Q5000 \\
\midrule
-0.05 & 0.250 & 667.6 & 1,775.4 & 3,139.6 & 5,512.3 \\
-0.01 & 0.290 & 667.5 & 1,774.3 & 3,147.4 & 5,403.6 \\
+0.00 & 0.300 & 667.5 & 1,774.5 & 3,148.0 & 5,403.0 \\
+0.01 & 0.310 & 667.3 & 1,774.6 & 3,148.5 & 5,381.3 \\
+0.05 & 0.350 & 666.8 & 1,770.3 & 3,148.9 & 5,303.3 \\
\bottomrule
\end{tabularx}
\begin{tablenotes}[flushleft]
\footnotesize
\item Entries are predicted income at t+1 (quetzales); column headings give income at t. General sample, labor income, population-weighted LOESS. -- indicates no supported prediction. Variant is the additive change from the selected span. Data and weights are held fixed.
\end{tablenotes}
\end{threeparttable}
\end{table}
